% ---
% id: receiver/detector-imperfections/req:efficiency-normalization
% kind: req
% title: "Compilation of a probe-free receiver at its normalized station"
% summary: >-
%   By default the engine compiles a probe-free receiver at the factor
%   t = η_max/(99/100), or t = 1 where η_max ≥ 99/100, and emits that station's
%   blocks, weights and dark-composed family. A declared factor in [η_max, 1]
%   replaces it, and one outside is refused.
% prominence: standard
% status: implemented
% uses: [prop:efficiency-normalization, def:normalized-station, rem:normalization-gain, rem:normalization-scope]
% assumes: []
% refs: []
% provenance: newly-derived
% implements: certiqs_key.receiver.efficiency:normalize_efficiencies
% lean:
% user_statement: >-
%   Unless you declare otherwise, a receiver without a probe source is analysed
%   with nearly all of its largest detector efficiency moved into the channel
%   as loss: every efficiency is divided by t, the largest one divided by 99/100,
%   so the best detector is analysed at 99/100; at a largest efficiency of
%   99/100 or more none is divided. This leaves your observed statistics
%   unchanged and typically raises the certified key rate at low efficiencies.
%   You may declare the factor, between the largest efficiency and one;
%   declaring one turns the normalization off. Declaring the largest
%   efficiency itself is allowed and sound, but with dark counts at a photon
%   cutoff of two or more the engine can then report no key where your model
%   gives one, which is why, below 99/100, the default stops short of it; a
%   receiver whose largest efficiency is one or close to it can meet the same
%   limitation.
% notes: >-
%   Where eta_max < 99/100 the default stops short of the full normalization
%   t = eta_max by the exact factor 99/100; where eta_max >= 99/100 it is
%   t = 1, which at eta_max = 1 is the full normalization. At t = eta_max a station with dark counts at photon cutoff 2
%   or more can have a singular two-photon block evaluated in floating point,
%   which the snap lifts and whose lifted directions the support check counts,
%   so a run can report key 0 where the model's rate is positive
%   (the receiver's design.md, structural decision 14); the margin may
%   change once the engine handles a singular block at the full
%   normalization.
%   The margin is tested at photon cutoffs 1 to 3 under the engine's snap and
%   support threshold. It moves the rate, never the bound:
%   prop:efficiency-normalization holds at every t in [eta_max, 1].
% ---

A receiver's measurement is uniform loss followed by its \tref{normalized station}
(Proposition~\ref{prop:efficiency-normalization}), and the weight a
consumer reads loses a factor of the squared efficiency without it
(Remark~\ref{rem:normalization-gain}); the factor itself is a choice, of which
Remark~\ref{rem:normalization-scope} gives the trade-off. The reason the
default stops short of the full normalization where the largest efficiency
is below $99/100$ is the engine's present handling of a singular block, not a
property of the identity, and is stated in
\repodoc{spec/modules/receiver/design.md}.

\begin{requirement}\label{req:efficiency-normalization}
For a request whose receiver declares no source modes and some efficiency
$\eta_k>0$, the engine compiles, in place of the declared receiver, its
normalized station at the \tref{factored loss} $t$, with the efficiency of port $k$
equal to $\eta_k/t$, and emits that station's blocks, multiclick weights and
\tref{dark-composed family}; $t$ is the value of \iface{efficiency_normalization}
where the request declares one; otherwise, with $\eta_{\max}$ the largest
declared efficiency, $t=\eta_{\max}/(99/100)$, formed exactly, where
$\eta_{\max}<99/100$, and $t=1$ where $\eta_{\max}\ge 99/100$. On every route the station compiled is the one at
the exact rational efficiencies $\eta_k/t$: its exact statements, its
reference and its enclosures are formed from those rationals, its
floating-point evaluation evaluates them as the \tref{evaluation model}
evaluates every declared number, and no value rounded from $\eta_k/t$ is
declared in their place; the quotient is exactly one at a port with
$\eta_k=t$. A declared $t$ outside
$[\eta_{\max},1]$ is refused with
\iface{efficiency_normalization_out_of_range}. A request that declares source
modes, or whose efficiencies are all zero, is compiled at $t=1$, and a
declared $t<1$ in such a request is refused with the same reason.
\end{requirement}
